\scanpage{79}
\mspage{7}
dont $l$ est justement un élément prenant le rôle d'un ``élément origine'', pour en déduire des $1$-cocycles\ldots

\deleted{Le théorème de Mordell-Weil implique que si $K$ est un corps de type fini sur $Q$, et si $X$ peut se plonger dans une \inserted{VA, ou dans un} groupe algébrique ext. d'une \uline{VA par un tore,} \uline{alors pour $x_0$ fixé,}}

Si on garde $x,y$, et qu'on remplace $l$ par un $l':x\longrightarrow y$, alors \deleted{\unsure{on}} on aura
\[
l'=la=bl\qquad a\in\Pi_1(X,x),\ b\in\Pi_1(X,y)
\]
et on aura, en vertu de (10)(11)
\[
\tag{31}\left\{\begin{aligned}
f_{l'}(u)&=\operatorname{int}(a^{-1})f_l(u)\underbrace{f_a(u)}_{a^{-1}\cdot{}^ua}=a^{-1}\cdot f_l(u)\cdot{}^ua\\
g_{l'}(u)&=\underbrace{g_b(u)}_{{}^ub\cdot b^{-1}}\operatorname{int}(b)\underbrace{g_l(u)}_{\deleted{\unsure{a^{-1}\cdot{}^ua}}}={}^ub\,g_l(u)b^{-1}
\end{aligned}\right.
\]
\[
(\text{i.e. }g_{l'}^{-1}(u)=b\,g_l^{-1}(u)\,{}^ub^{-1})
\]
\unsure{enfin}, $f_{l'}$ est un $1$-cocycle \inserted{à droite} homologue à $f_l$, $g_{l'}$ un $1$-cocycle \unsure{à g.} homologue à $g_l$. Ainsi, pour $x,y\in X_K$, on trouve des éléments
\[
\tag{32}\left\{\begin{aligned}
g_{y,x}&\in H^1_{\unsure{g}}(\Gamma,\widehat\Pi_1(X,y))\\
f_{y,x}&\in H^1_{\mathrm{dr}}(\Gamma,\widehat\Pi_1(X,x))
\end{aligned}\right.
\]
\deleted{\unsure{donc}} qui sont ``associés'' en un sens convenable, via les isom extérieurs canoniques
\diagram{\begin{tikzcd}[ampersand replacement=\&]
\widehat\Pi_1(X,x)\arrow[r,"\substack{\mathrm{isom}\\\mathrm{ext}}","\sim"']\&\widehat\Pi_1(\unsure{Y|X},y)
\end{tikzcd}}

\margin{Les $f_l(u),g_l(u)$ sont définis, et les formules \unsure{déjà de type|de type} (31) valables, sans \unsure{qu'on exige que} $l:x\longrightarrow y$ provienne de $\Pi_1(X)$, seulement de $\widehat\Pi_1(X)$. Et la définition des classes de $f_{\unsure{y}},g_{\unsure{y}}$ (ou classes de cohomologie) (32) ne dépend pas des (32) considérations transcendantes, ni de la \unsure{fixation du} corps de base $\mathbb C$\ldots}
\margin{78}
